\documentclass[11pt]{article}

\usepackage[a4paper, margin=1in]{geometry}
\usepackage[T1]{fontenc}
\usepackage[utf8]{inputenc}
\usepackage{amsmath, amssymb}
\usepackage{booktabs}
\usepackage{graphicx}
\usepackage{hyperref}
\usepackage{xcolor}
\hypersetup{colorlinks=true, linkcolor=blue, urlcolor=blue, citecolor=blue}
\usepackage{authblk}
\usepackage{microtype}
\usepackage{enumitem}
\setlist[itemize]{leftmargin=1.4em}
\setlist[enumerate]{leftmargin=1.6em}

\newcommand{\grad}{\nabla}

\title{\textbf{The Fine-Structure Constant as a Size Ratio}\\[0.3em]
\large H$^-$ and the Deuteron as Mirror Proton--Electron Models}

\author[1]{Claes Johnson}
\affil[1]{KTH Royal Institute of Technology, Stockholm, Sweden \\ \texttt{claesjohnson@gmail.com}}
\date{\today \\[0.4em]
\small\textit{Mathematical theory and formulation by the author; manuscript developed} \\
\small\textit{in collaboration with Claude (Anthropic). See Section~\ref{sec:collab}.}}

\begin{document}
\maketitle

\begin{abstract}
RealQM~\cite{realqm} and RealNucleus~\cite{realnucleus} model atoms and nuclei from the \emph{same}
proton--electron Coulomb mechanics --- charge densities minimising the Coulomb energy, no strong force, no
exchange--correlation functional. Their simplest mirror pair is the subject here: the \textbf{H$^-$ ion}
(one proton surrounded by two electrons in two shells) as the \emph{atom} model, and the \textbf{deuteron}
$^2$H (one electron caged inside two protons in two shells) as the \emph{nucleus} model --- the two related
by $p\leftrightarrow e$. Both models take only three inputs: the charge $\pm e$, an electron kinetic
coefficient $\kappa$, and the proton size $R_p$ --- and \textbf{no mass $m$, no speed of light $c$, no
Planck constant $\hbar$} (these never appear separately). The atomic size is $a_0=2\kappa/e^2$ (the
electrons taper to zero in vacuum, kinetic-set); the nuclear size is $R_p$ (the caged electron is flat, zero
kinetic energy, size-set). Their ratio is the fine-structure constant squared,
\[ \alpha^2 = \frac{R_p}{a_0},\qquad a_0 = \text{size of the H atom}, \]
so $\alpha=\sqrt{R_p/a_0}$ is a pure ratio of two lengths, independent of $m,c,\hbar$. With $R_p$ at the
classical-radius scale ($\approx2.8$ fm) the relation is exact; at the measured deuteron radius ($2.14$ fm)
it reproduces $\alpha$ to ${\sim}25\%$ (Section~\ref{sec:comp}). The two models are independently validated:
atomic total energies $<1\%$ across the periodic table, and the He-4/deuteron binding-energy ratio $13.1$
vs.\ $12.7$.
\end{abstract}

\section{Two Coulomb models: RealQM and RealNucleus}\label{sec:models}
The program treats protons and electrons as $\pm$ unit charges (charge densities) interacting by the Coulomb
law, with a kinetic penalty on the electron densities; structure is found by minimising the total energy.
\begin{itemize}
\item \textbf{RealQM}~\cite{realqm} --- \emph{atoms and molecules}: electron charge clouds arranged in
non-overlapping shells around small proton nuclei, bound by attraction to the nuclear charge. No basis set,
no exchange--correlation functional. It reproduces atomic total energies to $<1\%$ across the periodic
table~\cite{atoms}.
\item \textbf{RealNucleus}~\cite{realnucleus} --- \emph{nuclei}: protons bound by shared electrons (neutron
$=$ p$+$e), by the same Coulomb attraction, with \emph{no strong force}. Nested charge-cloud shells give,
e.g., the He-4/deuteron binding-energy ratio $13.1$ vs.\ $12.7$ experimental.
\end{itemize}
Both are the same mechanics --- charges and a kinetic coefficient --- run at two scales. The point of this
note is that the two scales are locked together by a single geometric relation, $\alpha^2=R_p/a_0$.

\section{The mirror pair: H$^-$ and the deuteron}\label{sec:mirror}
Take equal numbers of protons and electrons as $\pm$ charges on the same footing (the starting point of the
program; motivated cosmologically in~\cite{realcos}). The simplest three-body bound systems are exact
mirrors under $p\leftrightarrow e$:
\begin{center}
\begin{tabular}{lll}
\toprule
 & atom model (RealQM) & nucleus model (RealNucleus) \\
\midrule
system & \textbf{H$^-$} $=$ e\,+\,p\,+\,e & \textbf{deuteron} $^2$H $=$ p\,+\,e\,+\,p \\
picture & 1 proton, 2 electrons in 2 shells & 1 electron, 2 protons in 2 shells \\
centre / surround & proton centred; electrons around & electron caged; protons around \\
size scale & atomic, ${\sim}a_0$ & nuclear, ${\sim}R_p$ \\
\bottomrule
\end{tabular}
\end{center}
They are the \emph{same} 3-body Coulomb system --- a central charge holding two like-charged partners against
their mutual repulsion --- differing in nothing but which particle is central. The surrounders set the size:
electrons (fluid) $\Rightarrow$ atomic; protons (rigid, of fixed size $R_p$) $\Rightarrow$ nuclear.

\paragraph{The inputs of both models.} Charge $\pm e$; an electron kinetic coefficient $\kappa$ (taken as a
single primitive, not decomposed as $\hbar^2/2m$); and the proton size $R_p$. \textbf{No mass, no $c$, no
$\hbar$} enters separately --- only the combination $\kappa$, the charge $e$, and the length $R_p$. The
construction is non-relativistic and electrostatic.

\section{The two sizes, and $\alpha^2=R_p/a_0$}\label{sec:sizes}
An electron density $\psi$ ($\int\psi^2=1$) of size $L$ has energy
$E=\kappa\!\int|\grad\psi|^2 - Z\,e^2/L$.

\paragraph{Atomic size (charge-set).} Around a proton, an electron must taper to zero in the surrounding
vacuum; that varying profile costs kinetic energy ${\sim}\kappa/L^2$, and minimising $\kappa/L^2-e^2/L$ gives
\[ a_0=\frac{2\kappa}{e^2}\quad(=\hbar^2/me^2),\qquad\text{the size of the H atom.} \]

\paragraph{Nuclear size (size-set).} Caged between protons on its own domain with a multiphase \emph{free
boundary}, the electron's density does not drop to zero but hands off to the proton density; a flat caged
density has $\grad\psi=0$, hence \emph{zero kinetic energy}. With no kinetic pressure to spread it, its size
is set by the proton's finite size $R_p$ (the cage floor), binding Coulomb at that size.

\paragraph{The ratio.} Dividing the two sizes,
\begin{equation}
\boxed{\ \alpha^2=\frac{R_p}{a_0}=\frac{R_p\,e^2}{2\kappa},\qquad \alpha=\sqrt{R_p/a_0}.\ }
\end{equation}
Written this way $\alpha$ is a \textbf{ratio of two lengths}, obtained by measuring two sizes and dividing:
it carries \emph{no $m$, no $c$, no $\hbar$}. This is the textbook relation $r_e/a_0=\alpha^2$ (the classical
radius $r_e=e^2/mc^2$ over the Bohr radius) with $R_p$ identified as $r_e$; the standard form
$\alpha=e^2/\hbar c$ writes the same number through $c$ (and $m,\hbar$), but the ratio itself needs none of
them. In standard QM/QED $\alpha$ is the relativistic coupling (amplitude $\sqrt\alpha$ per photon vertex;
fine-structure splitting $\propto\alpha^2$; $v_1/c=\alpha$); here it is the mass-, $c$- and $\hbar$-free
proton-to-atom size ratio --- the more primitive reading.

\section{Computations}\label{sec:comp}
\paragraph{Atomic scale.} With the electron $\kappa$ and charge $e$,
\[ a_0=\frac{2\kappa}{e^2}=0.529\ \text{\AA}=52{,}918\ \text{fm}, \]
and RealQM reproduces the H atom at this size with total energy $-0.500$ Ha (exact).

\paragraph{Nuclear scale and the ratio.} Taking $R_p$ at candidate nuclear lengths:
\begin{center}
\begin{tabular}{lcccc}
\toprule
$R_p$ & value (fm) & $R_p/a_0$ & $\%$ of $\alpha^2$ & $\sqrt{\;}\to\alpha$ \\
\midrule
classical radius $r_e=e^2/mc^2$ & $2.818$ & $5.33\times10^{-5}$ & $100\%$ & $1/137$ (exact) \\
deuteron rms radius            & $2.142$ & $4.05\times10^{-5}$ & $76\%$  & $1/157$ \\
proton charge radius           & $0.841$ & $1.59\times10^{-5}$ & $30\%$  & $1/251$ \\
\bottomrule
\end{tabular}
\end{center}
(True $\alpha^2=5.33\times10^{-5}$, $\alpha=1/137.04$.) So with $R_p$ at the classical-radius scale the
relation $\alpha^2=R_p/a_0$ is exact by construction; at the \emph{measured deuteron radius} it reproduces
the fine-structure constant to ${\sim}25\%$ ($\alpha=1/157$ vs.\ $1/137$) --- from charge, $\kappa$ and the
proton size alone, with no strong force and no $c$. (The deuteron is unusually loosely bound and slightly
larger than a nucleon, so this is the \emph{scale/ratio} being right, not a precise fit.)

\paragraph{Independent validation of the two models.} The mirror models are quantitative on each side:
\begin{itemize}
\item \emph{Atom side} --- RealQM atomic \textbf{total energies $<1\%$} across the whole periodic table (all
$Z\ge11$; period-2 p-block B--F the only $3$--$6\%$ exception)~\cite{atoms}.
\item \emph{Nucleus side} --- RealNucleus \textbf{He-4/deuteron binding-energy ratio $13.1$ vs.\ $12.7$}
experimental (${\sim}3\%$, nested charge-cloud shells)~\cite{realnucleus}.
\end{itemize}
So $\alpha^2=R_p/a_0$ is bracketed by two independent quantitative validations --- the atomic scale to
$<1\%$ and the nuclear binding ratio to ${\sim}3\%$ --- beyond the ${\sim}25\%$ size check itself.

\section{Mass is a separate quantity}\label{sec:mass}
Nothing above used mass: H$^-$ and the deuteron are built from charge and $\kappa$. Mass is \textbf{inertia}
--- the localisation (frequency / inverse size) of the matter wave --- a \emph{separate, gravitational}
quantity, coupling to the electrostatic sector only through energy, $E=mc^2$. Electromagnetism (charge) is an
interaction; mass is the inverse size of each matter wave.

\section{Complexity from the size asymmetry, and open problems}\label{sec:open}
The proton--electron asymmetry that builds structure is not one of \emph{charge} (charge is exactly balanced,
$\pm e$) but of \emph{size}: the proton is fixed-and-small (a rigid anchor), the electron variable. The fixed
proton anchors; the variable electron spans scales --- shells in the atom, a cage-filling cloud in the
deuteron --- the nested hierarchy that \emph{is} complexity. And $\alpha$ measures this asymmetry:
$\alpha^2=R_p/a_0$ is the electron's inverse dynamic range ($1/\alpha^2\approx18{,}800$), so the smallness of
$\alpha$ is the engine of complexity.

\medskip\noindent\emph{Open.} (i) The nuclear \textbf{primitive}: $R_p$ as input length, or equivalently the
classical radius $r_e=\alpha^2a_0$ (electromagnetic mass~\cite{mmr}); $\alpha$'s value is reinterpreted, not
derived. (ii) Whether the solver realises a flat, zero-kinetic-energy caged electron at the proton scale is
not cleanly confirmed. (iii) The equal-footing treatment of $p$ and $e$ vs.\ the measured mass ratio.

\section{Speculative outlook}\label{sec:outlook}
The same ingredients extend, far more speculatively, to a cosmology --- equal-mass $\pm$ charges from two
fluctuating potentials, a gravitational bounce with negative-mass dark energy, and a primordial helium
fraction $Y\approx0.245$ --- in a companion essay~\cite{realcos}; noted here only for scope.

\section{Collaboration}\label{sec:collab}
The mathematical theory and the proposal are the author's. The manuscript was developed in dialogue with
Claude (Anthropic), which assisted in articulating the argument and its computations. We present the
human--AI collaboration transparently as part of the working method.

\begin{thebibliography}{9}
\bibitem{realqm} C.\ Johnson, \emph{RealQM: Quantum Mechanics as Multiphase 3D Continuum Mechanics},
manuscript and interactive gallery, \url{https://claes542.github.io/RealMolecule/}.
\bibitem{realnucleus} C.\ Johnson, \emph{RealNucleus: nuclear binding from a pure-Coulomb proton--electron
model} (neutron $=$ p$+$e; He-4/deuteron binding ratio from nested charge-cloud shells), manuscript, same
repository.
\bibitem{atoms} C.\ Johnson, \emph{RealQM total energies and ionization energies across the periodic table}
(atomic total energies within $1\%$ for $Z\ge11$), manuscript, same repository.
\bibitem{mmr} C.\ Johnson, \emph{Many-Minds Relativity} --- light as a medium wave, mass as field
self-energy ($P=mc\Rightarrow E=mc^2$), and a critique of Einstein's special relativity.
\bibitem{realcos} C.\ Johnson, \emph{Real Cosmology: a proton--electron universe from two fluctuating
potentials} (companion essay), same repository.
\end{thebibliography}

\end{document}
